%============================================================
% CHƯƠNG 1: GIỚI THIỆU ĐỀ TÀI
%============================================================
\chapter{GIỚI THIỆU ĐỀ TÀI}
\label{chap:intro}

Chương này trình bày tổng quan về đồ án, làm rõ bối cảnh, động cơ, mục
tiêu, phạm vi nghiên cứu và những đóng góp dự kiến. Xu hướng tích hợp
trí tuệ nhân tạo vào hệ thống IoT (AIoT) và nhu cầu cấp thiết về giám
sát, dự báo chất lượng không khí theo thời gian thực tạo ra cơ hội cho
việc phát triển các giải pháp quan trắc môi trường thông minh với chi
phí thấp. Đồ án hướng đến việc xây dựng một hệ thống hoàn chỉnh gồm
các node cảm biến IoT, pipeline xử lý dữ liệu chuỗi thời gian, mô hình
dự báo LSTM và giao diện trực quan hóa phân bố không gian, chứng minh
tính khả thi của AIoT trong lĩnh vực quan trắc môi trường đô thị.

%------------------------------------------------------------
\section{Đặt vấn đề}
\label{sec:intro:problem}
%------------------------------------------------------------

Ô nhiễm không khí hiện là một trong những thách thức môi trường và y tế
nghiêm trọng nhất của thế kỷ 21. Theo Tổ chức Y tế Thế giới (WHO), mỗi
năm có khoảng 7 triệu người tử vong sớm do tiếp xúc với không khí ô
nhiễm~\cite{who2021air}. Tại Việt Nam, tình trạng ô nhiễm không khí tại
các thành phố lớn như Hà Nội và Thành phố Hồ Chí Minh ngày càng trở
nên nghiêm trọng, với chỉ số AQI thường xuyên vượt ngưỡng an toàn, đặc
biệt trong các tháng mùa đông và những ngày không có gió~\cite{vea2023}.

Nhiệt độ và độ ẩm không khí là hai thông số môi trường cơ bản nhưng có
mối tương quan chặt chẽ với chất lượng không khí và sức khỏe con
người~\cite{who2021air}. Độ ẩm cao kết hợp với nhiệt độ thích hợp tạo
điều kiện cho sự phát triển của vi khuẩn, nấm mốc và tăng cường phản
ứng hóa học của các chất ô nhiễm trong không khí. Việc giám sát liên
tục, dự báo sớm xu hướng biến đổi của hai chỉ số này tại nhiều điểm
phân tán trong không gian là nền tảng thiết yếu cho các hệ thống cảnh
báo môi trường thông minh.

Tuy nhiên, các hệ thống quan trắc môi trường hiện có đang bộc lộ nhiều
hạn chế đáng kể:

\begin{itemize}[leftmargin=2cm]
  \item \textbf{Chi phí cao và mật độ thấp:} Các trạm quan trắc tự động
    của nhà nước có giá lên đến hàng trăm triệu đồng mỗi trạm, dẫn đến
    mật độ trạm thưa thớt, không đủ để phản ánh sự phân bố môi trường
    theo không gian trong các tòa nhà hay khu vực quy mô nhỏ.

  \item \textbf{Thiếu khả năng dự báo:} Phần lớn các hệ thống chỉ dừng
    lại ở việc đo và hiển thị dữ liệu hiện tại, không có khả năng dự
    báo xu hướng trong tương lai để người dùng chủ động ứng phó.

  \item \textbf{Thiếu trực quan hóa phân bố không gian:} Các hệ thống
    thường hiển thị dữ liệu dạng số hoặc biểu đồ thời gian, chưa cung
    cấp bản đồ nhiệt (heatmap) thể hiện sự phân bố môi trường theo
    không gian thực, gây khó khăn cho việc xác định vùng cần can thiệp.

  \item \textbf{Hệ sinh thái đóng:} Nhiều sản phẩm thương mại không hỗ
    trợ tùy biến hoặc mở rộng, gây khó khăn cho nghiên cứu và tích hợp
    vào các hệ thống hiện có.
\end{itemize}

Sự phát triển của \textbf{Internet of Things (IoT)} cho phép triển khai
mạng lưới node cảm biến phân tán chi phí thấp với khả năng kết nối và
truyền dữ liệu liên tục. Kết hợp với các kỹ thuật \textbf{học máy}
(Machine Learning), đặc biệt là các mô hình học sâu cho dữ liệu chuỗi
thời gian như LSTM (Long Short-Term Memory), hệ thống có thể không chỉ
quan trắc mà còn dự báo xu hướng môi trường trong tương lai. Hơn nữa,
các thuật toán nội suy không gian như \textbf{Kriging} và \textbf{IDW}
(Inverse Distance Weighting) cho phép tái tạo bản đồ phân bố môi trường
liên tục từ một tập hợp hữu hạn các điểm đo, mở ra khả năng trực quan
hóa trực quan và toàn diện hơn.

%------------------------------------------------------------
\section{Mục tiêu và phạm vi đề tài}
\label{sec:intro:scope}
%------------------------------------------------------------

\subsection{Mục tiêu đề tài}
\label{subsec:intro:objective}

Đồ án đặt ra mục tiêu xây dựng một \textbf{hệ thống AIoT hoàn chỉnh}
có khả năng giám sát và dự báo chất lượng không khí (tập trung vào
nhiệt độ, độ ẩm và AQI) trong một không gian xác định, bao gồm các mục tiêu
cụ thể sau:

\begin{enumerate}[leftmargin=2cm]
  \item \textbf{Xây dựng node cảm biến IoT:} Thiết kế và triển khai các
    node cảm biến sử dụng vi điều khiển ESP32 tích hợp cảm biến DHT11
    để thu thập dữ liệu nhiệt độ, độ ẩm và AQI mô phỏng định kỳ. Dữ liệu được truyền
    về hệ thống trung tâm qua giao thức MQTT qua kết nối Wi-Fi.

  \item \textbf{Xây dựng hệ thống backend thu thập và lưu trữ:} Phát
    triển backend Node.js/Express tích hợp MQTT broker (Aedes) để tiếp
    nhận dữ liệu từ các node cảm biến, lưu trữ vào cơ sở dữ liệu
    PostgreSQL và phân phối theo thời gian thực đến các client thông qua
    Socket.io.

  \item \textbf{Triển khai mô hình dự báo chuỗi thời gian:} Xây dựng
    và huấn luyện mô hình LSTM bằng PyTorch để dự báo nhiệt độ, độ
    ẩm và AQI trong 15 phút tới dựa trên 30 mẫu dữ liệu gần nhất. Mô hình
    được phục vụ thông qua một Python/FastAPI service độc lập.

  \item \textbf{Nội suy và trực quan hóa phân bố không gian:} Từ dữ
    liệu thực đo và dữ liệu dự báo của các node, áp dụng thuật toán nội
    suy Kriging và IDW để tạo lưới dữ liệu liên tục (40$\times$50 điểm),
    trực quan hóa dưới dạng heatmap trên bản đồ Leaflet.

  \item \textbf{Xây dựng giao diện người dùng:} Phát triển dashboard web
    React để hiển thị dữ liệu thời gian thực qua biểu đồ Chart.js, kết
    quả dự báo, heatmap phân bố không gian và hệ thống cảnh báo khi các
    chỉ số vượt ngưỡng.

  \item \textbf{Chi phí thấp và mã nguồn mở:} Sử dụng linh kiện phổ
    biến, giá rẻ và công bố mã nguồn mở để cộng đồng nghiên cứu có thể
    tái sử dụng và phát triển.
\end{enumerate}

\subsection{Phạm vi đề tài}
\label{subsec:intro:limitation}

Đề tài phát triển một sản phẩm prototype để chứng minh tính khả thi của
các mục tiêu đề ra. Phạm vi nghiên cứu được giới hạn cụ thể như sau:

\begin{itemize}[leftmargin=2cm]
  \item \textbf{Phần cứng:} Sử dụng vi điều khiển ESP32 kết hợp cảm
    biến DHT11 để đo nhiệt độ, độ ẩm và sinh AQI mô phỏng. Đề tài không thiết kế PCB tùy
    chỉnh và không mở rộng sang các cảm biến khí khác (PM2.5, CO$_2$).

  \item \textbf{Mô hình ML:} Tập trung vào kiến trúc LSTM (PyTorch) cho
    bài toán dự báo đa biến (nhiệt độ, độ ẩm) với window size 30 và
    forecast horizon 15 phút. Mô hình chạy phía server (FastAPI), không
    triển khai trực tiếp trên vi điều khiển.

  \item \textbf{Nội suy không gian:} Áp dụng hai phương pháp IDW và
    Kriging (thư viện \texttt{pykrige}, \texttt{scipy}) cho không gian
    phòng/khu vực. Không thực hiện mô hình khí tượng hoặc CFD.

  \item \textbf{Dữ liệu:} Mô hình được huấn luyện trên dữ liệu thu
    thập từ hệ thống prototype trong môi trường thực nghiệm. Không sử
    dụng dữ liệu khí tượng quốc gia.

  \item \textbf{Triển khai:} Hệ thống được thử nghiệm trong môi trường
    phòng lab với một số node cảm biến. Không thực hiện triển khai quy
    mô lớn ngoài trời hay đánh giá độ chính xác theo tiêu chuẩn đo
    lường quốc gia.
\end{itemize}

%------------------------------------------------------------
\section{Đóng góp của đồ án}
\label{sec:intro:contribution}
%------------------------------------------------------------

\subsection{Xây dựng bản thiết kế tổng thể hệ thống AIoT quan trắc
môi trường}
\label{subsec:intro:contrib1}

Đồ án đề xuất và xây dựng một bản thiết kế hoàn chỉnh cho hệ thống
giám sát và dự báo môi trường dựa trên AIoT theo kiến trúc bốn tầng:
tầng thu thập (node ESP32/DHT11), tầng truyền dẫn (MQTT/Wi-Fi), tầng
xử lý và lưu trữ (Node.js/Express, PostgreSQL, FastAPI/LSTM) và tầng
hiển thị (React dashboard). Thiết kế được xây dựng theo hướng mô-đun
hóa, sử dụng công nghệ phổ biến và chi phí thấp, cho phép tái sử dụng
và mở rộng dễ dàng cho các ứng dụng quan trắc môi trường khác.

\subsection{Triển khai pipeline dự báo chuỗi thời gian tích hợp
nội suy không gian}
\label{subsec:intro:contrib2}

Đồ án xây dựng thành công pipeline hoàn chỉnh từ huấn luyện mô hình
LSTM (PyTorch) trên dữ liệu chuỗi thời gian thực, phục vụ inference
qua FastAPI, đến kết hợp kết quả dự báo với thuật toán nội suy không
gian Kriging/IDW để tạo ra heatmap phân bố môi trường dự báo. Đây là
điểm kỹ thuật nổi bật, kết hợp dự báo thời gian và mô hình hóa không
gian trong cùng một hệ thống thời gian thực.

\subsection{Triển khai thành công hệ thống prototype hoạt động đầy đủ}
\label{subsec:intro:contrib3}

Dựa trên bản thiết kế đề xuất, đồ án triển khai thành công một prototype
hoạt động đầy đủ, bao gồm: các node cảm biến ESP32 giao tiếp qua MQTT,
backend xử lý và phân phối dữ liệu thời gian thực, Python service chạy
mô hình LSTM và tính toán heatmap, và dashboard web trực quan. Prototype
vận hành ổn định trong điều kiện thử nghiệm, chứng minh tính khả thi
của toàn bộ pipeline từ thu thập dữ liệu đến dự báo và trực quan hóa
không gian.

%------------------------------------------------------------
\section{Bố cục đồ án}
\label{sec:intro:structure}
%------------------------------------------------------------

Phần còn lại của báo cáo đồ án tốt nghiệp này được tổ chức như sau:

\textbf{Chương~\ref{chap:survey}} đi vào khảo sát yêu cầu thực tế của
bài toán thông qua phân tích các hệ thống quan trắc môi trường hiện có.
Trên cơ sở khảo sát, các chức năng cốt lõi của hệ thống được xác định
và biểu diễn qua biểu đồ use case kèm đặc tả chức năng chi tiết.

\textbf{Chương~\ref{chap:tech}} giới thiệu các công nghệ và công cụ
được lựa chọn sử dụng trong đồ án, bao gồm: vi điều khiển ESP32 và cảm
biến DHT11, giao thức MQTT, framework Node.js/Express, PostgreSQL, mô
hình LSTM (PyTorch), FastAPI, các thuật toán nội suy không gian
Kriging/IDW, và React với Leaflet.

\textbf{Chương~\ref{chap:design}} trình bày chi tiết quá trình thiết
kế và triển khai hệ thống: kiến trúc tổng thể bốn tầng, thiết kế
firmware ESP32, sơ đồ cơ sở dữ liệu PostgreSQL, quy trình huấn luyện
và triển khai mô hình LSTM, thiết kế FastAPI service, và giao diện
dashboard. Chương cũng bao gồm kết quả kiểm thử và thông tin triển
khai.

\textbf{Chương~\ref{chap:contribution}} tập trung làm nổi bật ba giải
pháp kỹ thuật đặc sắc: pipeline dự báo chuỗi thời gian LSTM phục vụ
qua FastAPI, cơ chế nội suy không gian Kriging/IDW để tạo heatmap thời
gian thực, và kiến trúc truyền thông MQTT kết hợp Socket.io đảm bảo độ
trễ thấp cho toàn hệ thống.

\textbf{Chương~\ref{chap:conclusion}} tổng kết quá trình thực hiện đề
tài, đánh giá kết quả đạt được so với mục tiêu ban đầu và đề xuất các
hướng phát triển trong tương lai, bao gồm mở rộng cảm biến, cải thiện
độ chính xác mô hình dự báo và nâng cấp hạ tầng triển khai.
